% !TeX root = ../../Book2.tex
\section{二进分解}
\label{b2:sec:1903}

势用一个齐次核同时处理所有频率。二进分解则把频率空间拆成比例相同的环带，每带对应一种尺度；它适合记录平滑阶数，并允许不同尺度以平方而非绝对和的方式合成。

\subsection{频率局部化}
\label{b2:subsec:190301}

取实值、径向、光滑函数 \(\chi\)，满足 \(0\leq\chi\leq1\)，在 \(|\xi|\leq1\) 时为一，在 \(|\xi|\geq2\) 时为零，且沿半径单调。令 \(\varphi(\xi)=\chi(\xi)-\chi(2\xi)\)。它支撑于 \(1/2\leq|\xi|\leq2\)，并满足
\[
 \sum_{j\in\mathbb Z}\varphi(2^{-j}\xi)=1\quad(\xi\ne0).
\]
等式由有限和的望远镜消去得到，不能在零频同样成立。频率定位的空间代价是失去紧支集：\(\check\varphi\) 是快速衰减函数，但一般遍布全空间。

\ProseParagraphBreak

一个位置紧支集而又频率紧支集的非零光滑函数不存在。Fourier 积分在紧频带时延拓为整函数，若它在位置上的开区间或开球外为零，唯一性迫使恒为零。后章的波包因此使用快速衰减的空间局部化，而不要求两种支集同时精确紧致。

\subsection{Littlewood--Paley 投影}
\label{b2:subsec:190302}

定义\term[littlewood-paley-touying]{Littlewood--Paley 投影}[Littlewood--Paley projection]
\[
 \Delta_jf\coloneqq\mathcal F^{-1}
 \bigl(\varphi(2^{-j}\,\cdot)\widehat f\bigr)
 =2^{jd}\check\varphi(2^j\,\cdot)*f.
\]
\symidx{\(\Delta_jf\)：Littlewood--Paley 频率投影}
这里的 \(\Delta_j\) 是频率投影记号，不是 Laplace 算子，也不是第一卷的有限差分 \(\Delta_h\)。下标的整数或向量类型区别其含义。
\begin{proposition}\label{b2:thm:19-lp-resolution}
若 \(1<p<\infty\)，则对 \(f\in L^p(\mathbb R^d)\)，有
\[
 f=\lim_{N\to\infty}\sum_{j=-N}^N\Delta_jf
 \quad\text{在 }L^p\text{ 中}.
\]
\end{proposition}
\begin{proof}
有限和乘子为 \(\chi(2^{-N}\xi)-\chi(2^{N+1}\xi)\)。第一项的空间核是积分一的 Schwartz 逼近恒等，故范数趋于 \(f\)。第二项是极低频平均。先对 \(f\in L^1\cap L^p\)，Young 不等式给出其 \(L^p\) 范数不超过 \(C2^{-(N+1)d(1-1/p)}\|f\|_1\)，趋于零。其卷积核的 \(L^1\) 范数一致有界，再以稠密性推广到一般 \(f\in L^p\)。在 \(p=1\) 时指数为零，不能使用这个结论。
\end{proof}

\subsection{二进环带}
\label{b2:subsec:190303}

每个 \(\Delta_jf\) 的频率支撑在半径约 \(2^j\) 的环带，相邻环带只有有限重叠。若 \(|j-k|>2\)，则 \(\Delta_j\Delta_k=0\)。投影并不是两两正交，也一般不是幂等；名称表示光滑频率定位，不能把它们当成 Hilbert 空间的正交投影。
下面采用\term[bernstein-budengshi]{Bernstein 不等式}[Bernstein inequality]的下述版本。

\begin{proposition}[Bernstein 估计]\label{b2:thm:19-bernstein}
若 \(f\in L^p\)、\(1\leq p\leq q\leq\infty\)，且 Fourier 支撑包含于 \(B(0,C2^j)\)，则对非负整数 \(k\)，
\[
 \|D^kf\|_q\leq C_{k,p,q}2^{j(k+d(1/p-1/q))}\|f\|_p.
\]
\end{proposition}
\begin{proof}
取光滑紧支集乘子 \(\eta\)，在该频带上为一。于是 \(D^kf=D^k(2^{jd}\check\eta(2^j\,\cdot))*f\)。令 \(1+1/q=1/p+1/r\)，核的 \(L^r\) 范数由换元等于 \(C2^{j(k+d(1-1/r))}\)。Young 不等式给出公式。恒等式先按分布成立，右侧卷积是所需范数中的函数，故也给出导数的函数代表。
\end{proof}
只知道 Fourier 支集半径还不能得到导数的反向界，因为低频函数可能落在球中心；反向估计需要环带避开零频。

% 扩写定位：b2:subsec:190303
环带避开原点时，一阶导数确实能反向控制函数。设 Fourier 支撑在
\(\{c2^j\leq|\xi|\leq C2^j\}\)，取光滑 \(\eta\) 在此环带上为一且支集仍避开零频。令
\[
 a_\ell(\zeta)=\frac{\eta(\zeta)\zeta_\ell}{2\pi\mathrm i|\zeta|^2}.
\]
则 \(\sum_\ell a_\ell(2^{-j}\xi)\,(2\pi\mathrm i\xi_\ell)=2^j\) 于该支集，故
\[
 f=2^{-j}\sum_{\ell=1}^d a_\ell(2^{-j}D)\partial_\ell f,\qquad
 \|f\|_p\leq C2^{-j}\sum_\ell\|\partial_\ell f\|_p
 \quad(1\leq p\leq\infty).
\]
每个 \(a_\ell\) 光滑紧支集，其逆 Fourier 核可积，缩放保持核 \(L^1\) 范数，因而只用 Young 即可证明这一界。这里并未提前调用后面的 Mihlin 定理；环带上的紧频率符号与全空间零次齐次符号有不同的核可积性。

\ProseParagraphBreak

把这一步重复 \(k\) 次，或者直接以多指标配分 \(|\xi|^{2k}\)，得到
\[
 \|f\|_p\leq C_k2^{-jk}\sum_{|\alpha|=k}\|\partial^\alpha f\|_p.
\]
它与正向 Bernstein 共同说明：一个纯环带函数的一阶导数，尺度就是 \(2^j\)；频率越高，相同函数范数所需的梯度越大。一个低频常数则不能满足这样的反向不等式，因为它的梯度为零。非齐次 Sobolev 范数中的低频 \(L^p\) 项正是为了补回这一信息。

\subsection{近似正交性}
\label{b2:subsec:190304}

在 \(L^2\) 中，由 Plancherel，
\[
 \sum_j\|\Delta_jf\|_2^2
 =\int\left(\sum_j|\varphi(2^{-j}\xi)|^2\right)|\widehat f(\xi)|^2\,\mathrm d\xi.
\]
分割非负且和为一，重叠次数不超过固定整数 \(L\)，故括号在非零频处介于 \(1/L\) 与一之间。这给出平方和与 \(L^2\) 范数的等价，常数依赖所选分割而不是输入。

\ProseParagraphBreak

若一列函数分别具有相互分离的频率支集，则其平方范数可以相加；若支集只有限重叠，就在重叠次数内损失一个常数。这是\term[ji-hu-zhengjiao]{几乎正交性}[almost orthogonality]的最基本版本。对 \(L^p\) 不能直接使用 Plancherel，下一节用随机符号与奇异积分把同样的尺度信息传递到全部开区间指数。
